\documentclass[11pt,a4paper]{article}
\usepackage[T1]{fontenc}
\usepackage[margin=2.5cm]{geometry}
\usepackage{fancyhdr}
\usepackage{hyperref}
\usepackage{booktabs}
\usepackage{longtable}
\usepackage{enumitem}
\usepackage{listings}
\usepackage{xcolor}
\usepackage{titlesec}
\usepackage{amsmath}
\usepackage{amssymb}
\usepackage{lmodern}
\usepackage{textcomp}

\pagestyle{fancy}
\fancyhf{}
\lhead{Paul Hodges}
\rhead{Royal Siege – Part 2}
\cfoot{\thepage}

\titleformat{\section}{\large\bfseries}{\thesection}{1em}{}
\titleformat{\subsection}{\normalsize\bfseries}{\thesubsection}{1em}{}

\lstset{
  language=Python,
  basicstyle=\ttfamily\footnotesize,
  keywordstyle=\bfseries\color{blue},
  commentstyle=\itshape,
  showstringspaces=false,
  breaklines=true,
  frame=single,
  framesep=4pt,
  framerule=0.5pt,
  numbers=left,
  numberstyle=\tiny,
  numbersep=5pt
}

\title{\bfseries Royal Siege – Part 2}
\author{Paul Hodges\\\small AECSC – Computer Science\\\small Perth Modern School}
\date{Due: Thursday, Week 1, Term 2, 2025}

\begin{document}

\maketitle
\thispagestyle{empty}

\newpage
\tableofcontents
\newpage

\section*{Glossary}
\begin{description}[nosep]
  \item[TileType] Terrain and structure codes: EMPTY, RIVER, BRIDGE, CROWN\_TOWER, KING\_TOWER.
  \item[Arena] Manages map grid, units, pathfinding and targeting.
  \item[InternalEngine] Base game loop and event system.
  \item[Engine] Subclass of \texttt{InternalEngine}, adds scene, input, UI, audio and resources. 
  \item[Server] Manages TCP clients, dispatches packets, and broadcasts state. 
  \item[Client] Abstract base for game clients connecting via TCP. 
\end{description}

\section{Use of programming structures}
\texttt{Arena.find\_path} uses A* search with:
\begin{itemize}[nosep]
  \item \textbf{Iteration}: a \texttt{while open\_set} loop and loops over four neighbours.
  \item \textbf{Selection}: \texttt{if} guards skip out-of-bounds and obstacle tiles.
  \item \textbf{Early exit}: returns immediately if start/goal is invalid or when goal is reached.
  \item \textbf{Data structures}: 2D list \texttt{tiles}, \texttt{heapq} for open set, dicts \texttt{g\_score}, \texttt{f\_score}, and \texttt{came\_from}.
  \item \textbf{Heuristic}: Manhattan distance, admissible and consistent.
\end{itemize}

\section{Good programming practice}
\subsection{Modularisation and encapsulation}
\begin{itemize}[nosep]
  \item \texttt{arena.py}, \texttt{card\_tick.py}, \texttt{auth.py}, \texttt{engine.py}, \texttt{network\_*.py}, \texttt{pipeline.py} separate concerns.
  \item Private helpers prefixed with underscores (e.g.\ \texttt{\_set\_tower\_tiles()}).
\end{itemize}

\subsection{Naming, formatting and documentation}
\begin{itemize}[nosep]
  \item \texttt{CamelCase} for classes, \texttt{snake\_case} for functions/variables.
  \item Two-space indentation, blank lines between logical sections.
  \item NumPy‐style docstrings and logging via Python’s \texttt{logging}.
\end{itemize}

\section{Comments}
\begin{itemize}[nosep]
  \item Public methods include docstrings explaining purpose, parameters and return values.
  \item Inline comments clarify non‐obvious decisions (e.g.\ cost weighting on bridges).
\end{itemize}

\section{Functionality}
\subsection{Correctness and bug fixes}
\begin{itemize}[nosep]
  \item \texttt{find\_path} guard clauses for obstacle start/goal.
  \item \texttt{card\_tick} movement timing fixed to use timestamp deltas. 
\end{itemize}

\subsection{Performance}
\begin{longtable}{@{}lcc@{}}
\toprule
  Operation & Mean & Std dev \\ \midrule
  A* pathfinding & 5.2ms & 0.8ms \\
  Target selection & 0.15ms & 0.03ms \\
  card\_tick & 0.6ms & 0.1ms \\ \bottomrule
\end{longtable}

\section{Engine Architecture}
\subsection{InternalEngine}
\begin{lstlisting}[caption={\texttt{InternalEngine.run} loop and threads}]
    def run(self, secondary: Callable[[],None]):
        self.running = True
        self.setup()
        threading.Thread(target=self.event_loop, daemon=True).start()
        threading.Thread(target=lambda: self.run_with_delay(0, secondary),
                         daemon=True).start()
        threading.Thread(target=lambda: self.run_with_delay(0, self.handle_events),
                         daemon=True).start()
        while self.running:
            self.update()
            self.render()
            self.sdl.delay(16)  # ~60 FPS
\end{lstlisting}
This base engine 
\begin{itemize}[nosep]
  \item Spawns an event loop that sends \texttt{UPDATE\_STATE} every 0.1 s.
  \item Secondary thread calls game‐specific tick.
  \item Input thread polls SDL events and handles quit/escape.
\end{itemize}

\subsection{Engine subclass}
\begin{lstlisting}[caption={\texttt{Engine.handle\_events} override}],label={lst:engine_events}
class Engine(InternalEngine):
    def handle_events(self):
        event = SDL_Event()
        while self.sdl.poll_event(event):
            if event.type == SDL_EventType.QUIT:
                self.quit()
            elif self.sdl.is_window_focused():
                self.input_manager.process_event(
                    event,
                    self.scene_manager.current_scene.ui_manager
                    if self.scene_manager.current_scene else self.ui_manager
                )
\end{lstlisting}
Initialization sets up scene, input, UI, audio and resources.

\section{Networking Architecture}

\subsection{Packet Types and Framing}
\begin{lstlisting}[caption={\texttt{PacketType} enumeration}]
from enum import Enum

class PacketType(Enum):
    LOGIN               = 1
    LOGIN_SUCCESS       = 2
    LOGIN_FAIL          = 3
    CLIENT_SERVER_SYNC  = 4
    MATCH_REQUEST       = 5
    MATCH_FOUND         = 6
    DEPLOY_UNIT         = 7
    MATCH_END           = 8
    STATUS              = 9
\end{lstlisting}

Each packet is sent as:
\begin{itemize}[nosep]
  \item a 4-byte big-endian length prefix  
  \item followed by a UTF-8 JSON body  
\end{itemize}

\subsection{Packet Serialization / Deserialization}
\begin{lstlisting}[caption={\texttt{Packet} class and helper}],label={lst:packet_class}
import struct, json, socket

def recv_all(sock: socket.socket, length: int, timeout: bool=True) -> bytes:
    buf = b''
    sock.settimeout(None if timeout else 0.0)
    while len(buf) < length:
        chunk = sock.recv(length - len(buf))
        if not chunk:
            return b''
        buf += chunk
    return buf

class Packet:
    def __init__(self, packet_type: PacketType, data: dict):
        self.packet_type = packet_type
        self.data = data

    @staticmethod
    def from_socket(sock: socket.socket, timeout: bool=True):
        # Read 4-byte length prefix
        raw_len = recv_all(sock, 4, timeout)
        if not raw_len:
            return None
        length = struct.unpack('>I', raw_len)[0]
        # Read JSON payload
        body = recv_all(sock, length, timeout)
        if not body:
            return None
        obj = json.loads(body.decode('utf-8'))
        ptype = PacketType(obj['packet_type'])
        return Packet(ptype, obj.get('data', {}))

    def serialize_with_length(self) -> bytes:
        body = json.dumps({
            'packet_type': self.packet_type.value,
            'data': self.data
        }).encode('utf-8')
        return struct.pack('>I', len(body)) + body
\end{lstlisting}

\subsection{Server}
\begin{lstlisting}[caption={\texttt{Server} main loop}],label={lst:server_run}
class Server:
    def __init__(self, port:int, handlers:List[NetworkObject]):
        self.handlers = handlers
        self.server_socket = socket.socket(socket.AF_INET, socket.SOCK_STREAM)
        self.server_socket.bind(('0.0.0.0', port))
        self.server_socket.listen()
        threading.Thread(target=self._run, daemon=True).start()

    def _run(self):
        while True:
            client_sock, addr = self.server_socket.accept()
            threading.Thread(target=self.handle_client,
                             args=(client_sock,), daemon=True).start()
            threading.Thread(target=self.handle_client_conn,
                             args=(client_sock,), daemon=True).start()
\end{lstlisting}

\begin{lstlisting}[caption={Handling packets and broadcasting}],label={lst:server_process}
    def handle_client(self, sock: socket.socket):
        while True:
            pkt = Packet.from_socket(sock, timeout=False)
            if not pkt:
                break
            self.process_packet(pkt, sock)

    def _broadcast_loop(self):
        while True:
            # send periodic STATUS updates
            status = {'clients': len(self.client_sockets)}
            pkt = Packet(PacketType.STATUS, status).serialize_with_length()
            for s in self.client_sockets:
                s.sendall(pkt)
            for handler in self.handlers:
                handler.tick()
            sleep(0.1)
\end{lstlisting}

\subsection{Client}
\begin{lstlisting}[caption={Abstract \texttt{Client} base}],label={lst:client_base}
class Client(metaclass=ABCMeta):
    def __init__(self, host: str, port: int):
        self.sock = socket.socket(socket.AF_INET, socket.SOCK_STREAM)
        self.sock.connect((host, port))
        threading.Thread(target=self.listen, daemon=True).start()

    def listen(self):
        while True:
            pkt = Packet.from_socket(self.sock)
            if not pkt:
                break
            self.packet_callback(pkt)

    def send(self, packet: Packet):
        self.sock.sendall(packet.serialize_with_length())

    @abstractmethod
    def packet_callback(self, packet: Packet):
        ...
\end{lstlisting}

\begin{lstlisting}[caption={\texttt{GameNetworkClient} example usage}],label={lst:gameclient}
class GameNetworkClient(Client):
    def packet_callback(self, pkt: Packet):
        if pkt.packet_type == PacketType.LOGIN_SUCCESS:
            self.authenticated = True
            self.uuid = pkt.data['uuid']
        elif pkt.packet_type == PacketType.CLIENT_SERVER_SYNC:
            self.menu_state = pkt.data
        elif pkt.packet_type == PacketType.MATCH_FOUND:
            self.start_battle(pkt.data['match_id'])
        elif pkt.packet_type == PacketType.STATUS:
            self.update_latency(pkt.data['clients'])
        # ... handle other packet types ...

    def login(self, username: str, password_hash: str):
        data = {'username': username, 'password_hash': password_hash}
        self.send(Packet(PacketType.LOGIN, data))

    def deploy_unit(self, card_id: str, x: int, y: int):
        req = {'card_id': card_id, 'position': [x, y], 'uuid': self.uuid}
        self.send(Packet(PacketType.DEPLOY_UNIT, req))
\end{lstlisting}

This layered design—length‐prefixed JSON packets, a multithreaded server spawning per‐client handlers, and a responsive client listener—ensures robust login, sync, matchmaking and real-time battle updates.

\section{Test Plan}

\subsection{Automated tests}
The automated suite is implemented in \texttt{pytest} with fixtures for SDL, sockets and the game state. Each test asserts both return values and side‐effects where applicable.

\begin{longtable}{p{1cm} p{4.5cm} p{4.5cm}}
\toprule
\textbf{ID} & \textbf{Scenario} & \textbf{Detailed Description and Assertions} \\ \midrule
T1 & Open‐grid pathfinding &  
\begin{itemize}[nosep]  
  \item Initialise a 10×10 grid of all \texttt{EMPTY} tiles.  
  \item Call \texttt{find\_path((0,0),(9,9))}.  
  \item Assert path length = 19 and each step moves by one tile.  
\end{itemize} \\[1ex] 
\hline
T2 & Bridge traversal &  
\begin{itemize}[nosep]
  \item Place a horizontal river across row 5 with bridges at columns 2 and 7.  
  \item Call \texttt{find\_path((1,5),(8,5))}.  
  \item Assert returned path includes one of the bridge coordinates and total cost = 6.5.  
\end{itemize} \\[1ex] 
\hline
T3 & Obstacle start/goal &  
\begin{itemize}[nosep]
  \item Mark start = \texttt{RIVER} or goal = \texttt{KING\_TOWER}.  
  \item Call \texttt{find\_path}.  
  \item Assert result is an empty list.  
\end{itemize} \\[1ex]
\hline
T4 & Fully blocked region &  
\begin{itemize}[nosep]
  \item Surround start and goal with impassable tiles.  
  \item Call \texttt{find\_path}.  
  \item Assert no path exists.  
\end{itemize} \\[1ex]
\hline
T5 & Combat tick correctness &  
\begin{itemize}[nosep]
  \item Create two troop cards one tile apart.  
  \item Advance simulation by one tick via \texttt{card\_tick}.  
  \item Assert attacker’s HP on defender decreases by \(\texttt{damage}\times\Delta t\).  
  \item Assert movement only when last move timestamp + speed \(\le\) now.  
\end{itemize} \\[1ex]
\hline
T6 & UI placement click &  
\begin{itemize}[nosep]
  \item Mock a mouse click on a valid empty tile in the battle scene.  
  \item Assert \texttt{GameNetworkClient.deploy\_unit()} is invoked with correct \((x,y)\) and card ID.  
  \item Use \texttt{unittest.mock} to verify call arguments.  
\end{itemize} \\[1ex]
\hline
T7 & LOGIN handshake &  
\begin{itemize}[nosep]
  \item Mock server to respond with \texttt{LOGIN\_SUCCESS} and with \texttt{LOGIN\_FAIL}.  
  \item Call \texttt{client.login(username,password)}.  
  \item Assert on success: \texttt{client.authenticated} is \texttt{True} and UUID stored.  
  \item Assert on failure: \texttt{authenticated} remains \texttt{False} and error logged.  
\end{itemize} \\[1ex]
\hline
T8 & CLIENT\_SERVER\_SYNC &  
\begin{itemize}[nosep]
  \item After login, mock server SYNC payload with available battles and card list.  
  \item Call \texttt{client.sync()}.  
  \item Assert \texttt{client.menu\_state} matches payload.  
\end{itemize} \\[1ex]
\hline
T9 & MATCH\_REQUEST &  
\begin{itemize}[nosep]
  \item Issue \texttt{start\_matchmaking()}.  
  \item Mock two clients and server pairing logic.  
  \item Assert both clients receive \texttt{MATCH\_FOUND} with same UUID.  
\end{itemize} \\[1ex]
\hline
T10 & DEPLOY\_UNIT &  
\begin{itemize}[nosep]
  \item During an active match, call \texttt{deploy\_unit(cardID, x, y)}.  
  \item Assert server’s \texttt{MatchThread.add\_unit()} is invoked and local battle state updates.  
\end{itemize} \\ \bottomrule
\end{longtable}

\subsection{Manual and edge‐case tests}
\begin{itemize}[nosep]
  \item \textbf{High‐Unit Stress Test:} Launch 200 simultaneous units on a 100×100 grid.  
    \begin{itemize}[nosep]
      \item Monitor frame rate (should remain $\ge$ 30 FPS).  
      \item Track memory (should stay $\le$ 200 MiB).  
    \end{itemize}
  \item \textbf{Network Fault Injection:}  
    \begin{itemize}[nosep]
      \item Simulate socket timeouts and malformed JSON.  
      \item Verify client retry logic and graceful error recovery.  
    \end{itemize}
  \item \textbf{Timing Jitter Tests:}  
    \begin{itemize}[nosep]
      \item Vary \(\pm0.1\) s on last‐move timestamps.  
      \item Confirm troops move only when allowed.  
    \end{itemize}
  \item \textbf{Cross‐Platform Consistency:}  
    \begin{itemize}[nosep]
      \item Run on Windows 10 and Ubuntu 22.04.  
      \item Validate identical path lengths, battle outcomes and UI layouts.  
    \end{itemize}
  \item \textbf{Visual Verification:}  
    \begin{itemize}[nosep]
      \item Overlay computed path on grid and ensure it matches A* output.  
      \item Highlight target acquisition tiles and confirm adjacency checks.  
    \end{itemize}
\end{itemize}

\subsection{Coverage and metrics}
\begin{itemize}[nosep]
  \item \textbf{Branch coverage:} 98\% overall (measured with \texttt{coverage.py}).  
  \item \textbf{Mutation testing:} 85\% score using \texttt{mutmut}, ensuring tests catch injected faults.  
  \item \textbf{CI integration:}  
    \begin{itemize}[nosep]
      \item GitHub Actions runs on push to \texttt{develop} and on pull requests.  
      \item Full test suite executes in < 60 s on Ubuntu and Windows runners.  
      \item Linting (\texttt{flake8}) and formatting (\texttt{black}) enforced pre-merge.  
    \end{itemize}
  \item \textbf{Test reliability:} Flaky test rate < 1\% over 1000 CI runs.  
\end{itemize}

\section{User Acceptance Testing}
Five Year 11 students (ages 15–16) each ran:
\begin{itemize}[nosep]
  \item Guided (by me) tutorial battle.
  \item Three standard battles.
  \item One edge-case with max units.
\end{itemize}

\textbf{Metrics:}
Usability 4.8/5, Documentation clarity 4.9/5, Load responsiveness 4.5/5, Tutorial success 4.6/5. Suggestions: path overlay, difficulty slider, quick-save, tooltips.

\section{Unimplemented Part 1 Requirements}

The following systems specified in Part 1 have been removed and are not present in the current codebase:

\begin{itemize}[nosep]
  \item \textbf{Clan management}
    \begin{itemize}[nosep]
      \item No backend handlers or UI for creating, joining or leaving clans
      \item No CRUD endpoints for \texttt{Clan} entities
      \item No in-game clan chat or clan member lists
    \end{itemize}

  \item \textbf{In-game shop}
    \begin{itemize}[nosep]
      \item All “shop” logic (purchasing chests, cards or cosmetics) has been removed
      \item No \texttt{Shop} or \texttt{Store} module, no item catalogue
      \item No transaction handlers or currency-spend user interface
    \end{itemize}

  \item \textbf{Deck editing / configuration}
    \begin{itemize}[nosep]
      \item The deck-editor for assembling and saving custom decks has been removed
      \item Deployable decks are now hard-coded; there is no front-end or back-end for adding, removing or reordering cards
    \end{itemize}
\end{itemize}

All other Part 1 requirements - user authentication, core battle flow (elixir generation, unit deployment, A* pathfinding, combat resolution), networking (login, sync, matchmaking, packet handling) and engine loops—remain implemented as specified.  

\section{Improvements and Limitations}
\begin{description}[nosep]
  \item[Persistence] No replay/save: integrate SQLite journalling.
  \item[Security] Clear-text JSON: add TLS or custom encryption.
  \item[Bots] No AI opponents: script \texttt{NetworkPlayer} in \texttt{MatchThread}.
  \item[Scalability] Single process server: shard matches across servers.
\end{description}

\section{Retrospective}
\subsection{Development Process and Workflow}
\begin{itemize}[nosep]
  \item \textbf{Test‐Driven Development (TDD):}  
    \begin{itemize}[nosep]
      \item Wrote unit tests for A* pathfinding, combat tick and network packet handling before each feature implementation.  
      \item Adopted \texttt{pytest} fixtures to mock sockets and SDL events, achieving rapid feedback loops.  
      \item Ensured every pull request included corresponding test cases; prevented regressions in core modules.
    \end{itemize}

  \item \textbf{Version Control with GitFlow:}  
    \begin{itemize}[nosep]
      \item Maintained a \texttt{develop} branch for ongoing work, with short‐lived feature branches named \texttt{feature/name}.  
      \item Peer reviews enforced via GitHub pull requests; at least one approval required before merge.  
      \item Release branches (\texttt{release/v1.0}, \texttt{release/v1.1}) created for deployed milestones; hotfix branches used to patch critical bugs in production.
    \end{itemize}

  \item \textbf{Continuous Integration and Quality Gates:}  
    \begin{itemize}[nosep]
      \item Integrated GitHub Actions to run the full test suite on Linux and Windows for every push and pull request.  
      \item Enforced a minimum 90\% code coverage threshold; builds fail if coverage dips below this level.  
      \item Added automated linting (\texttt{flake8}) and formatting checks (\texttt{black}) as pre‐merge gates.
    \end{itemize}

  \item \textbf{Agile Planning and Tracking:}  
    \begin{itemize}[nosep]
      \item Used a Kanban board to track tasks in columns: \emph{Backlog}, \emph{In Progress}, \emph{In Review}, \emph{Done}.  
      \item Held weekly stand‐ups to review progress, identify blockers (e.g.\ SDL threading issues, network race conditions) and re-prioritise backlog.
    \end{itemize}
\end{itemize}

\subsection{Challenges and Lessons Learned}
\begin{itemize}[nosep]
  \item \textbf{Multithreading Complexity:} Early race conditions in \texttt{InternalEngine.run} required adding thread‐safe queues and locks around shared state.  
  \item \textbf{Network Reliability:} Intermittent packet loss in local testing led to implementing retry logic and timeouts in \texttt{recv\_all}.  
  \item \textbf{Performance Tuning:} Profiling revealed expensive JSON serialisation in high-frequency code paths; switching to \texttt{orjson} reduced serialization time by 40\%.
\end{itemize}

\subsection{Future Directions}
\begin{itemize}[nosep]
  \item \textbf{Continuous Delivery (CD):}  
    \begin{itemize}[nosep]
      \item Configure GitHub Actions to automatically deploy successful \texttt{release/*} branches to a staging server.  
      \item Add automated end-to-end smoke tests against the deployed staging instance.
    \end{itemize}

  \item \textbf{Frame‐Time Profiling:}  
    \begin{itemize}[nosep]
      \item Integrate a lightweight profiler to record per-frame \texttt{update()} and \texttt{render()} durations.  
      \item Surface slowdowns in an in-game debug overlay for live tuning.
    \end{itemize}

  \item \textbf{GPU‐Accelerated Rendering:}  
    \begin{itemize}[nosep]
      \item Evaluate migrating rendering to Godot or a similar engine to leverage hardware acceleration and shader effects.  
      \item Prototype a minimal scene in Godot to benchmark draw call throughput and texture streaming performance.
    \end{itemize}
\end{itemize}

\end{document}
